\subsection{Matrix Operations and Identities}
Throughout this subsection, all matrices are over the same fixed field $\mathbb{F}$, and every matrix size is a positive integer.

\begin{definition} \label{def:matrix-trace-and-determinant} \hlt{Trace and Determinant of a Matrix}\\
Let $n\geq1$ and $A=[a_{jk}]\in\operatorname{Mat}(n,n,\mathbb{F})$.
\begin{enumerate}[label=(\roman*)]
\setlength{\itemsep}{0pt}
\item The \hlt{trace} of $A$ is $\operatorname{tr}(A)\defeq\sum\limits_{j=1}^{n}a_{jj}$.
\item The \hlt{determinant} of $A$ is $\det A\defeq\sum\limits_{\sigma\in\mathfrak{S}_n}\sign(\sigma)\prod\limits_{j=1}^{n}a_{\sigma(j),j}$.
\end{enumerate}
\end{definition}

\begin{theorem} \label{thm:determinant-basic-laws} \hlt{Basic Determinant Laws}\\
Let $n\geq1$ and $A,B\in\operatorname{Mat}(n,n,\mathbb{F})$.
\begin{enumerate}[label=(\roman*)]
\setlength{\itemsep}{0pt}
\item $\det I_n=1$.
\item For fixed $k\in\{1,\ldots,n\}$ and $(A_{\cdot,j})_{j\neq k}$, the map $\mathbf{x}\mapsto\det(A_{\cdot,1},\ldots,\mathbf{x},\ldots,A_{\cdot,n})$ is linear.
\item Swapping two columns negates the determinant.
\item If two columns of $A$ are equal, then $\det A=0$.
\item If $A$ is triangular, then $\det A=\prod\limits_{j=1}^{n}a_{jj}$.
\item $\det(AB)=\det(A)\det(B)$.
\item $A$ is invertible if and only if $\det A\neq0$.
\item If $A$ is invertible, then $\det(A^{-1})=(\det A)^{-1}$.
\item $\det(A^t)=\det(A)$.
\end{enumerate}
\begin{proof}
\begin{enumerate}[label=(\roman*)]
\setlength{\itemsep}{0pt}
\item[$(i)$] In the permutation formula for $I_n$, only the identity permutation has a nonzero product, and its contribution is $1$.
\item[$(ii)$] For $k\in\{1,\ldots,n\}$, $\det(A_{\cdot,1},\ldots,\mathbf{x},\ldots,A_{\cdot,n})=\sum\limits_{\sigma\in\mathfrak{S}_n}\sign(\sigma)x_{\sigma(k)}\prod\limits_{\ell\neq k}a_{\sigma(\ell),\ell}$, which is linear in $\mathbf{x}$.
\item[$(iii)$] Let $\tau$ transpose columns $p$ and $q$, and let $B_{\cdot,j}=A_{\cdot,\tau(j)}$. Then $\det B=-\det A$ by reindexing the determinant sum with $\rho=\sigma\circ\tau$, because $\sign(\rho\circ\tau)=-\sign(\rho)$.
\item[$(iv)$] If columns $p$ and $q$ are equal, pair the term indexed by $\sigma$ with the term indexed by $\sigma\circ(p\ q)$. Their products agree and their signs are opposite, so every pair sums to zero, including when $\operatorname{char}(\mathbb{F})=2$.
\item[$(v)$] Suppose $A$ is upper triangular. A nonzero term satisfies $(\forall j\in\{1,\ldots,n\})\,\sigma(j)\leq j$, hence $\sigma=\operatorname{id}$ and $\det A=\prod\limits_{j=1}^{n}a_{jj}$. The lower-triangular case is analogous with $(\forall j\in\{1,\ldots,n\})\,\sigma(j)\geq j$.
\item[$(vi)$] Let $F$ be alternating and multilinear in its $n$ column arguments. Standard-basis expansion produces terms indexed by $(i_1,\ldots,i_n)\in\{1,\ldots,n\}^n$; repeated indices vanish by alternation.\\
The surviving tuples are permutations, and $F(\mathbf{e}_{\sigma(1)},\ldots,\mathbf{e}_{\sigma(n)})=\sign(\sigma)F(I_n)$ by alternation. Hence $F(A_{\cdot,1},\ldots,A_{\cdot,n})=F(I_n)\det A$.\\
For fixed $A$, $H(B)\defeq\det(AB)$ is alternating and multilinear in the columns of $B$, with $H(I_n)=\det A$. Thus $\det(AB)=H(B)=H(I_n)\det B=\det(A)\det(B)$.
\item[$(vii)$] ($\Rightarrow$) If $A$ is invertible, then $\det A\neq0$ because $1=\det(AA^{-1})=\det(A)\det(A^{-1})$ by parts (i) and (vi).
\item[$(vii)$] ($\Leftarrow$) If $A$ were singular, one column would be a linear combination of the others, and $\det A=0$ by parts (ii) and (iv), contrary to $\det A\neq0$. Thus $A$ is invertible.
\item[$(viii)$] If $A$ is invertible, then $\det(A^{-1})=(\det A)^{-1}$ by $1=\det(AA^{-1})=\det(A)\det(A^{-1})$ from parts (i), (vi), and (vii).
\item[$(ix)$] $\det(A^t)=\sum\limits_{\sigma\in\mathfrak{S}_n}\sign(\sigma)\prod\limits_{j=1}^{n}a_{j,\sigma(j)}=\sum\limits_{\tau\in\mathfrak{S}_n}\sign(\tau)\prod\limits_{j=1}^{n}a_{\tau(j),j}=\det A$ by reindexing with $\tau=\sigma^{-1}$.
\end{enumerate}
\hlt{Structural Commentary}: The permutation formula establishes alternation and triangularity before multiplicativity and invertibility.\\
\hlt{Key Insight}: Alternating multilinearity reduces every column expansion to signed basis permutations.
\end{proof}
\end{theorem}

\begin{lemma} \label{lem:determinant-identity-direct-sum} \hlt{Identity Blocks Do Not Change Determinant}\\
Let $m\geq1$, $n\geq0$, and $A\in\operatorname{Mat}(m,m,\mathbb{F})$.
\begin{enumerate}[label=(\roman*)]
\setlength{\itemsep}{0pt}
\item $\det(I_n\oplus A)=\det A$;
\item $\det(A\oplus I_n)=\det A$.
\end{enumerate}
\begin{proof}
\begin{enumerate}[label=(\roman*)]
\setlength{\itemsep}{0pt}
\item If $n=0$, the identity holds by the empty identity-block convention. If $n\geq1$, every nonzero permutation term fixes $\{1,\ldots,n\}$ pointwise; relabeling the remaining indices produces exactly the determinant expansion of $A$.
\item If $n=0$, the identity again holds by convention. If $n\geq1$, every nonzero term fixes the last $n$ indices pointwise, and the remaining sum is $\det A$.
\end{enumerate}
\hlt{Structural Commentary}: Each block position reduces the permutation sum to the indices occupied by $A$.\\
\hlt{Key Insight}: An identity block forces its coordinates to remain fixed.
\end{proof}
\end{lemma}

\begin{theorem} \label{thm:block-upper-triangular-determinant} \hlt{Determinant of an Upper Block-Triangular Matrix}\\
Let $r,s\geq1$, $A\in\operatorname{Mat}(r,r,\mathbb{F})$, $B\in\operatorname{Mat}(r,s,\mathbb{F})$, and $D\in\operatorname{Mat}(s,s,\mathbb{F})$.\\
Then $\det\left(\begin{smallmatrix}A&B\\0&D\end{smallmatrix}\right)=\det(A)\det(D)$.
\begin{proof}
$\left(\begin{smallmatrix}A&B\\0&D\end{smallmatrix}\right)=\left(\begin{smallmatrix}I_r&0\\0&D\end{smallmatrix}\right)\left(\begin{smallmatrix}I_r&B\\0&I_s\end{smallmatrix}\right)\left(\begin{smallmatrix}A&0\\0&I_s\end{smallmatrix}\right)$.\\
The outer determinants are $\det D$ and $\det A$ by Lemma \ref{lem:determinant-identity-direct-sum}; the middle determinant is $1$ by Theorem \ref{thm:determinant-basic-laws}(v).\\
The formula follows by Theorem \ref{thm:determinant-basic-laws}(vi).\\
\hlt{Structural Commentary}: A three-factor decomposition isolates the two diagonal-block determinants.\\
\hlt{Key Insight}: The off-diagonal block lies in a unit upper-triangular factor.
\end{proof}
\end{theorem}

\begin{definition} \label{def:polynomial-evaluation-on-square-matrix} \hlt{Polynomial Evaluation on a Square Matrix}\\
Let $n\geq1$ and $A\in\operatorname{Mat}(n,n,\mathbb{F})$. Define $A^0\defeq I_n$ and $(\forall\,j\in\mathbb{Z}_{\geq0})\,A^{j+1}\defeq A^jA$.\\
For $p(z)=\sum\limits_{j=0}^{d}c_jz^j\in\mathbb{F}[z]$, where $d\geq0$, the \hlt{evaluation of $p$ at $A$} is $p(A)\defeq\sum\limits_{j=0}^{d}c_jA^j\in\operatorname{Mat}(n,n,\mathbb{F})$.
\end{definition}

\begin{proposition} \label{prop:direct-sum-matrix-algebra} \hlt{Algebra of Matrix Direct Sums}
\begin{enumerate}[label=(\roman*)]
\setlength{\itemsep}{0pt}
\item If $A,C\in\operatorname{Mat}(m,n,\mathbb{F})$, $B,D\in\operatorname{Mat}(r,s,\mathbb{F})$, and $a,b\in\mathbb{F}$, then $a(A\oplus B)+b(C\oplus D)=(aA+bC)\oplus(aB+bD)$.
\item If $A\in\operatorname{Mat}(m,n,\mathbb{F})$, $C\in\operatorname{Mat}(n,\ell,\mathbb{F})$, $B\in\operatorname{Mat}(r,s,\mathbb{F})$, and $D\in\operatorname{Mat}(s,t,\mathbb{F})$, then $(A\oplus B)(C\oplus D)=AC\oplus BD$.
\item If $A\in\operatorname{Mat}(m,m,\mathbb{F})$, $B\in\operatorname{Mat}(n,n,\mathbb{F})$, and $k\in\mathbb{Z}_{>0}$, then $(A\oplus B)^k=A^k\oplus B^k$.
\item If $A\in\operatorname{Mat}(m,m,\mathbb{F})$, $B\in\operatorname{Mat}(n,n,\mathbb{F})$, and $p\in\mathbb{F}[z]$, then $p(A\oplus B)=p(A)\oplus p(B)$.
\end{enumerate}
\begin{proof}
\begin{enumerate}[label=(\roman*)]
\setlength{\itemsep}{0pt}
\item By Definition \ref{def:matrix-direct-sum}, Definition \ref{def:matrices-and-matrix-operations}, and Proposition \ref{prop:block-matrix-arithmetic}, $a(A\oplus B)+b(C\oplus D)=\left[\begin{smallmatrix}aA+bC&0\\0&aB+bD\end{smallmatrix}\right]=(aA+bC)\oplus(aB+bD)$.
\item By Definition \ref{def:matrix-direct-sum} and Proposition \ref{prop:block-matrix-arithmetic}, $(A\oplus B)(C\oplus D)=\left[\begin{smallmatrix}AC&0\\0&BD\end{smallmatrix}\right]=AC\oplus BD$.
\item Set $P_k\defeq((A\oplus B)^k=A^k\oplus B^k)$. By Definition \ref{def:polynomial-evaluation-on-square-matrix} and Proposition \ref{prop:matrix-product-associativity-and-identities}, $P_1$ holds.\\
If $P_k$ holds, then $(A\oplus B)^{k+1}=(A\oplus B)^k(A\oplus B)=(A^k\oplus B^k)(A\oplus B)=A^{k+1}\oplus B^{k+1}$ by the recurrence defining powers, the induction hypothesis, and (ii). Therefore $(\forall\,k\in\mathbb{Z}_{>0})\,P_k$ by Theorem \ref{thm:mathematical-induction}.
\item Write $p(z)=\sum\limits_{j=0}^{d}c_jz^j$. The identity $(A\oplus B)^j=A^j\oplus B^j$ holds for $j=0$ by Definition \ref{def:polynomial-evaluation-on-square-matrix}, Definition \ref{def:matrix-direct-sum}, for $j\geq1$ by (iii).\\
By (i), scalar multiples, finite sums of these direct sums are componentwise. Hence $p(A\oplus B)=p(A)\oplus p(B)$ by Definition \ref{def:polynomial-evaluation-on-square-matrix}.
\end{enumerate}
\hlt{Structural Commentary}: Block formulas extend by induction and linearity.\\
\hlt{Key Insight}: Direct-sum blocks evolve independently under matrix algebra.
\end{proof}
\end{proposition}

\begin{proposition} \label{prop:diagonal-multiplication-scales-rows-columns} \hlt{Diagonal Multiplication Scales Rows and Columns}\\
Let $A=[a_{ij}]\in\operatorname{Mat}(m,n,\mathbb{F})$, let $\lambda_1,\ldots,\lambda_n,\mu_1,\ldots,\mu_m\in\mathbb{F}$, and set $\Lambda\defeq\operatorname{diag}(\lambda_1,\ldots,\lambda_n)$ and $M\defeq\operatorname{diag}(\mu_1,\ldots,\mu_m)$. Then:
\begin{enumerate}[label=(\roman*)]
\setlength{\itemsep}{0pt}
\item $A\Lambda=[\lambda_ja_{ij}]$; right multiplication by $\Lambda$ scales column $j$ by $\lambda_j$.
\item $MA=[\mu_ia_{ij}]$; left multiplication by $M$ scales row $i$ by $\mu_i$.
\end{enumerate}
\begin{proof}
\begin{enumerate}[label=(\roman*)]
\setlength{\itemsep}{0pt}
\item For admissible $i,j$, $(A\Lambda)_{ij}=\sum\limits_{r=1}^{n}a_{ir}\Lambda_{rj}=a_{ij}\lambda_j=\lambda_ja_{ij}$ by Definition \ref{def:matrices-and-matrix-operations} and commutativity in $\mathbb{F}$.
\item For admissible $i,j$, $(MA)_{ij}=\sum\limits_{r=1}^{m}M_{ir}a_{rj}=\mu_ia_{ij}$ by Definition \ref{def:matrices-and-matrix-operations}.
\end{enumerate}
\hlt{Structural Commentary}: Diagonal sparsity isolates one summand in each entry formula.\\
\hlt{Key Insight}: The multiplication side determines whether rows or columns are scaled.
\end{proof}
\end{proposition}

\begin{proposition} \label{prop:block-upper-triangular-inverse} \hlt{Inverse of an Upper Block-Triangular Matrix}\\
Let $Y\in\operatorname{Mat}(m,m,\mathbb{F})$ and $Z\in\operatorname{Mat}(n,n,\mathbb{F})$ be invertible, and let $X\in\operatorname{Mat}(m,n,\mathbb{F})$.\\
Then $\left[\begin{smallmatrix}Y&X\\0&Z\end{smallmatrix}\right]^{-1}=\left[\begin{smallmatrix}Y^{-1}&-Y^{-1}XZ^{-1}\\0&Z^{-1}\end{smallmatrix}\right]\in\operatorname{Mat}(m+n,m+n,\mathbb{F})$, where each lower-left zero block has size $n\times m$.
\begin{proof}
Set $M\defeq\left[\begin{smallmatrix}Y&X\\0&Z\end{smallmatrix}\right]$ and $N\defeq\left[\begin{smallmatrix}Y^{-1}&-Y^{-1}XZ^{-1}\\0&Z^{-1}\end{smallmatrix}\right]$.\\
By Proposition \ref{prop:block-matrix-arithmetic} and Proposition \ref{prop:matrix-product-associativity-and-identities}, $MN=\left[\begin{smallmatrix}I_m&-XZ^{-1}+XZ^{-1}\\0&I_n\end{smallmatrix}\right]=I_{m+n}$ and $NM=\left[\begin{smallmatrix}I_m&Y^{-1}X-Y^{-1}X\\0&I_n\end{smallmatrix}\right]=I_{m+n}$.\\
Thus $N=M^{-1}$ by Definition \ref{def:matrices-and-matrix-operations} and Lemma \ref{lem:two-sided-inverse-uniqueness}.\\
\hlt{Structural Commentary}: Two block products verify both inverse identities.\\
\hlt{Key Insight}: The off-diagonal correction cancels the coupling block.
\end{proof}
\end{proposition}

\begin{theorem} \label{thm:inverse-upper-triangular} \hlt{Inverse of an Upper Triangular Matrix}\\
Let $n\geq1$ and let $A=[a_{ij}]\in\operatorname{Mat}(n,n,\mathbb{F})$ be upper triangular with $(\forall\,i\in\{1,\ldots,n\})\,a_{ii}\neq0$. Then:
\begin{enumerate}[label=(\roman*)]
\setlength{\itemsep}{0pt}
\item $A$ is invertible.
\item $A^{-1}$ is upper triangular.
\item The diagonal of $A^{-1}$ is $(a_{11}^{-1},\ldots,a_{nn}^{-1})$.
\end{enumerate}
\begin{proof}
Let $P_n$ denote the conjunction of (i)--(iii) in size $n$ under the stated hypotheses.\\
For $n=1$, $A^{-1}=[a_{11}^{-1}]$ by $[a_{11}][a_{11}^{-1}]=I_1=[a_{11}^{-1}][a_{11}]$, Definition \ref{def:matrices-and-matrix-operations}, Lemma \ref{lem:two-sided-inverse-uniqueness}.\\
This inverse is upper triangular with diagonal $(a_{11}^{-1})$, so $P_1$ holds. Suppose $n>1$ and $P_{n-1}$ holds.\\
Write $A=\left[\begin{smallmatrix}B&\mathbf{x}\\0&a_{nn}\end{smallmatrix}\right]$, where $B\in\operatorname{Mat}(n-1,n-1,\mathbb{F})$ is upper triangular with nonzero diagonal and $\mathbf{x}\in\mathbb{F}^{n-1}$.
\begin{enumerate}[label=(\roman*)]
\setlength{\itemsep}{0pt}
\item The blocks $B$, $[a_{nn}]$ invertible by induction hypothesis, Definition \ref{def:matrices-and-matrix-operations}, respectively. Hence $A$ invertible by Proposition \ref{prop:block-upper-triangular-inverse}.
\item The matrix $B^{-1}$ is upper triangular by the induction hypothesis.\\
By Proposition \ref{prop:block-upper-triangular-inverse}, $A^{-1}=\left[\begin{smallmatrix}B^{-1}&-B^{-1}\mathbf{x}a_{nn}^{-1}\\0&a_{nn}^{-1}\end{smallmatrix}\right]$, which is upper triangular by Definition \ref{def:matrices-and-matrix-operations}.
\item By the induction hypothesis, the diagonal of $B^{-1}$ is $(a_{11}^{-1},\ldots,a_{n-1,n-1}^{-1})$.\\
By the block formula in (ii), the diagonal of $A^{-1}$ is $(a_{11}^{-1},\ldots,a_{n-1,n-1}^{-1},a_{nn}^{-1})$.
\end{enumerate}
Thus $P_n$ holds, so $(\forall\,n\in\mathbb{Z}_{>0})\,P_n$ by Theorem \ref{thm:mathematical-induction}.\\
\hlt{Structural Commentary}: Induction transfers all three claims through the block inverse formula.\\
\hlt{Key Insight}: Triangular inversion reduces dimension while preserving the zero pattern.
\end{proof}
\end{theorem}
